% Generated draft: requires mathematical and source review.
\documentclass[11pt]{article}
\usepackage[margin=25mm]{geometry}
\usepackage{fontspec}
\setmainfont{Latin Modern Roman}
\usepackage{amsmath,amssymb,mathtools,unicode-math}
\setmathfont{Latin Modern Math}
\usepackage{enumitem,needspace,xurl,hyperref}
\setlength{\parindent}{0pt}
\setlength{\parskip}{5pt}
\setlength{\emergencystretch}{3em}
\newcommand{\N}{\mathbb{N}}
\newcommand{\Z}{\mathbb{Z}}
\newcommand{\Q}{\mathbb{Q}}
\newcommand{\R}{\mathbb{R}}
\newcommand{\C}{\mathbb{C}}
\newcommand{\F}{\mathbb{F}}
\newcommand{\A}{\mathbb{A}}
\newcommand{\PP}{\mathbb{P}}
\newcommand{\E}{\mathbb{E}}
\newcommand{\eps}{\varepsilon}
\newcommand{\dd}{\,\mathrm{d}}
\begin{document}
\section*{UnsolvedMath ID: 800237}
\textbf{Source problem code:} catalogue entry
\subsection{Definitions and mathematical statement}
The supplied source asks: \emph{Which geometric properties of the domain do uniform estimates depend on?} Its parenthetical example states that \(L^2\) estimates depend on the diameter and the dimension.

To explain the terminology without supplying missing hypotheses, a uniform estimate can be described schematically as follows. Given a specified family \(\mathcal D\) of domains, a specified class \(\mathcal U(D)\) of admissible objects on each \(D\in\mathcal D\), and specified nonnegative quantities \(A_D(u)\) and \(B_D(u)\), a bound is uniform over this family if
\[
\exists C\in[0,\infty)\quad\forall D\in\mathcal D\quad\forall u\in\mathcal U(D),\qquad A_D(u)\le C B_D(u).
\]
Dependence on geometric data would instead permit a bound controlled by a function of explicitly specified geometric quantities of \(D\). This is only an explanation of the question's terminology, not an identification of the estimates intended by the source.

\textbf{Unresolved formulation.} The supplied record specifies neither the ambient space and admissible domains nor an operator, equation, boundary condition, admissible data, norms, or estimate to be bounded. It also does not identify the parameters with respect to which uniformity is required. Consequently, it does not determine a precise quantified mathematical problem. The several complex variables context alone does not supply these hypotheses. The parenthetical \(L^2\) example is retained as a source assertion; it does not specify which \(L^2\) estimate is meant, whether diameter and dimension are sufficient parameters, or what other hypotheses are required.
\subsection{Short English statement}
The question asks which geometric features of a domain control constants in uniform estimates. The supplied example mentions diameter and dimension for \(L^2\) estimates, but the estimates and admissible domains are unspecified, so a precise mathematical target remains unresolved.
\subsection{Sources}
\begin{itemize}
\item Local UnsolvedMath registry, permanent ID \(800237\), exact source problem code \texttt{catalogue entry}, title \emph{Several Complex Variables Problem List --- Problem 34}. The supplied registry labels the problem open; this status has not been independently verified.
\item Buzzard, Fornaess, Krantz, \emph{The Several Complex Variables Problem List} (1995), Problem 34. Supplied source URL: \url{https://arxiv.org/abs/math/9509202}. The registry attributes the problem to Buzzard, Fornaess, Krantz and gives proposed year 1995.
\item Supplied provenance records access on \texttt{2026-10-16} and extraction method \texttt{arxiv\_latex}. This is supplied dated provenance, not an independently verified source visit or literature assessment.
\end{itemize}
\end{document}
